% 并查集：元素1..6，按边序列 {1,2},{3,4},{1,3},{5,6},{4,6} 合并；按大小合并，大小相同时把第二个根接到第一个根下。
% find(4) 在 (c) 中经 4->3->1，路径压缩后 4 直接指向 1。箭头由子点指向父点。
\begin{tikzpicture}[font=\figurefont\fontsize{8}{10}\selectfont,>=Stealth,
  el/.style={circle,draw=FigureStroke,line width=.8pt,minimum size=5.8mm,inner sep=0pt,text=FigureInk,fill=white},
  rt/.style={el,fill=FigureBlueFill},
  par/.style={draw=FigureStroke,line width=.9pt,->},
  newp/.style={draw=FigureGreen,line width=1.3pt,->},
  pathp/.style={draw=FigureYellow,line width=1.3pt,->,dashed},
  sz/.style={font=\FigureNoteFont\fontsize{7.5}{9}\selectfont,text=FigureMuted},
  note/.style={font=\FigureNoteFont\fontsize{7.5}{9}\selectfont,text=FigureMuted,align=center}]
% (a)
\begin{scope}
  \node[rt] (a1) at (0,0) {1}; \node[el] (a2) at (0,-1) {2};
  \node[rt] (a3) at (0.75,0) {3}; \node[el] (a4) at (0.75,-1) {4};
  \node[rt] (a5) at (1.5,0) {5}; \node[rt] (a6) at (2.15,0) {6};
  \draw[par] (a2)--(a1); \draw[par] (a4)--(a3);
  \node[align=center] at (1.07,-2.75) {(a) union(1,2)，\\union(3,4)};
\end{scope}
% (b)
\begin{scope}[xshift=3.15cm]
  \node[rt] (b1) at (0.45,0) {1}; \node[el] (b2) at (0,-1) {2};
  \node[el] (b3) at (0.9,-1) {3}; \node[el] (b4) at (0.9,-2.0) {4};
  \draw[par] (b2)--(b1); \draw[newp] (b3)--(b1); \draw[par] (b4)--(b3);
  \node[sz,right] at (b1.east) {大小 4};
  \node at (0.45,-2.55) {(b) union(1,3)};
  \node[note] at (0.45,-3.05) {大小相同，\\3接到1下};
\end{scope}
% (c)
\begin{scope}[xshift=5.75cm]
  \node[rt] (c1) at (0.45,0) {1}; \node[el] (c2) at (0,-1) {2};
  \node[el] (c3) at (0.9,-1) {3}; \node[el] (c4) at (0.9,-2.0) {4};
  \draw[par] (c2)--(c1); \draw[par] (c3)--(c1); \draw[pathp] (c4)--(c3);
  \draw[newp] (c4) to[out=20,in=-10] (c1);
  \node at (0.45,-2.55) {(c) find(4)};
  \node[note] at (0.45,-3.05) {经$4\to3\to1$，\\压缩后4改指1};
\end{scope}
% (d)
\begin{scope}[xshift=8.3cm]
  \node[rt] (d1) at (1.0,0) {1};
  \node[el] (d2) at (0,-1) {2}; \node[el] (d3) at (0.67,-1) {3}; \node[el] (d4) at (1.34,-1) {4};
  \node[el] (d5) at (2.0,-1) {5}; \node[el] (d6) at (2.0,-2.0) {6};
  \draw[par] (d2)--(d1); \draw[par] (d3)--(d1); \draw[par] (d4)--(d1);
  \draw[newp] (d5)--(d1); \draw[par] (d6)--(d5);
  \node[sz,right] at (d1.east) {大小 6};
  \node[align=center] at (1.0,-2.75) {(d) union(5,6)，\\union(4,6)};
\end{scope}
\end{tikzpicture}
